\section{Communication-Aware Adaptive KLA Fusion}
\label{sec:method}

This section describes the adaptive fusion rule used by the distributed GA-LMB filter. The starting point is conservative geometric-average Kullback--Leibler averaging for neighboring posteriors with unknown cross-correlations \cite{Battistelli2014KLA,Hlinka2014ICI}. The method allocates online fusion weights from bounded reliability cues. Each node then passes these weights to its GA-LMB merger. Appendix~\ref{app:theory} gives variational and information-geometric interpretations of the separated branches.

\subsection{Baseline GA-LMB fusion}

At node $s$, let $\mathcal{N}_s$ denote the local communication neighborhood. Let $\pi_k^{(j)}$ be the measurement-updated LMB posterior density formed from the delivered stream associated with neighbor $j \in \mathcal{N}_s$, as described in Section~\ref{sec:problem-formulation}. The conservative fusion target is the weighted geometric average in \eqref{eq:ga-kla-objective}. The remaining task is to specify how the LMB parameter sets are moment-projected and how the spatial and existence weights are consumed by the merger.

After moment projection, the Gaussian approximation of the $i$th Bernoulli spatial density at node $j$ is
\begin{equation}
p_{k,i}^{(j)}(x) \approx \mathcal{N}\!\left(x;\nu_{k,i}^{(j)}, T_{k,i}^{(j)}\right).
\end{equation}
Using spatial fusion weights $\omega_{k,s}^{x,(j)}$, the fused spatial density is obtained in canonical form:
\begin{equation}
\label{eq:method-gaussian-canonical}
K_{k,i}^{(s)} = \sum_{j \in \mathcal{N}_s} \omega_{k,s}^{x,(j)} \left(T_{k,i}^{(j)}\right)^{-1},
\qquad
h_{k,i}^{(s)} = \sum_{j \in \mathcal{N}_s} \omega_{k,s}^{x,(j)} \left(T_{k,i}^{(j)}\right)^{-1}\nu_{k,i}^{(j)},
\end{equation}
\begin{equation}
\Sigma_{k,i}^{(s)} = \left(K_{k,i}^{(s)}\right)^{-1},
\qquad
\mu_{k,i}^{(s)} = \Sigma_{k,i}^{(s)} h_{k,i}^{(s)}.
\end{equation}
Using existence fusion weights $\omega_{k,s}^{r,(j)}$, the Bernoulli existence probability is fused by the corresponding weighted geometric rule:
\begin{equation}
\label{eq:method-existence-geometric}
r_{k,i}^{(s)} =
\frac{\eta_{k,i}^{(s)} \prod_{j \in \mathcal{N}_s} \left(r_{k,i}^{(j)}\right)^{\omega_{k,s}^{r,(j)}}}
{\eta_{k,i}^{(s)} \prod_{j \in \mathcal{N}_s} \left(r_{k,i}^{(j)}\right)^{\omega_{k,s}^{r,(j)}} +
\prod_{j \in \mathcal{N}_s} \left(1-r_{k,i}^{(j)}\right)^{\omega_{k,s}^{r,(j)}}},
\end{equation}
where $\eta_{k,i}^{(s)}$ is the Gaussian normalization term induced by the fused spatial density.

\subsection{Adaptive reliability factors}

The adaptive layer is organized as a reliability-allocation design. It first forms a shared active-neighbor backbone that identifies informative fusion inputs. The next subsection routes this backbone into branch-specific spatial and existence weights.

\paragraph{Shared backbone.}
For node $s$ and neighbor $j \in \mathcal{N}_s$, define the raw score
\begin{equation}
\tilde{\omega}_{k,s}^{(j)} =
m_k^{(j)}
\cdot
q_{\mathrm{cov},k}^{(j)}
\cdot
q_{\mathrm{link},k}^{(j)}
\cdot
q_{\mathrm{exist},k}^{(j)},
\end{equation}
where $m_k^{(j)} \in \{0,1\}$ is the availability mask, $q_{\mathrm{cov},k}^{(j)}$ is a covariance-quality score, $q_{\mathrm{link},k}^{(j)}$ is a realized link-quality score, and $q_{\mathrm{exist},k}^{(j)}$ is an existence-confidence score. The normalized weight is
\begin{equation}
\bar{\omega}_{k,s}^{(j)} =
\frac{\tilde{\omega}_{k,s}^{(j)}}{\sum_{u \in \mathcal{N}_s}\tilde{\omega}_{k,s}^{(u)}}.
\end{equation}
This factorization addresses the three failure modes most directly tied to distributed fusion: diffuse local posteriors, missing delivered information, and ambiguous Bernoulli support. Figure~\ref{fig:factorization} shows how this shared backbone feeds the decoupled spatial and existence-weight paths used by the two operating modes.

\noindent\textbf{Remark 1 (backbone rationale).}
The product form is a compact way to combine independent reliability cues on the simplex. In log form, it adds the utilities $\log q_{\mathrm{cov}}$, $\log q_{\mathrm{link}}$, and $\log q_{\mathrm{exist}}$, so a poor value in any cue cannot be completely hidden by a large value in another cue. The availability mask $m_k^{(j)}$ is a hard gate for an unavailable neighbor, whereas the positive floors used in some cue mappings are soft safeguards for uncertain but still usable posteriors. If all active raw scores vanish, the implementation falls back to the uniform distribution over active neighbors before applying any smoothing or final floor.

\paragraph{Posterior concentration and realized delivery.}

For node $j$, let $T_{k,i}^{(j)}$ be the moment-matched covariance of the $i$th Bernoulli component. The covariance-quality score is defined as
\begin{equation}
q_{\mathrm{cov},k}^{(j)}
=
\frac{1}{\epsilon + \frac{1}{M_k}\sum_{i=1}^{M_k}\operatorname{tr}\!\left(T_{k,i}^{(j)}\right)},
\end{equation}
where $\epsilon>0$ is a small numerical constant used to avoid division by zero. A node with smaller posterior spread therefore receives a larger score.

The link-quality term is derived from realized communication outcomes. Let $d_k^{(j)}$ and $\ell_k^{(j)}$ denote the counts of delivered and dropped measurements, respectively. The realized link-quality score is defined as
\begin{equation}
q_{\mathrm{link},k}^{(j)}
=
\frac{d_k^{(j)}}{d_k^{(j)} + \ell_k^{(j)}},
\end{equation}
which directly reflects how much of that node's measurement stream remains available after delivery constraints.

\noindent\textbf{Design rationale.}
The covariance cue is a local posterior-concentration proxy: a smaller average trace means that, after moment projection, the neighbor contributes a sharper Gaussian component to the canonical-form spatial average above. The link cue uses delivered and dropped measurements at the current time step, so the weights reflect the information available to the local fusion problem.

\paragraph{Existence decisiveness.}

The third cue is separated because Bernoulli existence uncertainty differs from spatial posterior spread. For the $i$th Bernoulli component at node $j$,
\begin{equation}
c_{k,i}^{(j)} = \left|2r_{k,i}^{(j)} - 1\right|.
\end{equation}
The per-node weighted existence confidence is then
\begin{equation}
\bar{c}_{k}^{(j)} =
\frac{\sum_{i=1}^{M_k} r_{k,i}^{(j)} c_{k,i}^{(j)}}
{\epsilon + \sum_{i=1}^{M_k} r_{k,i}^{(j)}}.
\end{equation}
The final existence-confidence score is mapped into a bounded positive factor:
\begin{equation}
q_{\mathrm{exist},k}^{(j)}
=
\lambda_{\min} + (1-\lambda_{\min})\left(\bar{c}_{k}^{(j)}\right)^{p_e},
\end{equation}
where $\lambda_{\min} \in (0,1]$ and $p_e > 0$ are tunable parameters.
Figure~\ref{fig:existence-confidence} visualizes this bounded Bernoulli-level mapping and representative node-level aggregation effects.

\noindent\textbf{Remark 2 (existence confidence).}
The mapping $|2r-1|$ assigns low confidence to an undecided Bernoulli component near $r=0.5$ and high confidence to decisive absence or presence. Weighting this certainty by $r_{k,i}^{(j)}$ makes the node-level score emphasize components that are likely to exist. The lower bound $\lambda_{\min}$ weakens the influence of an uncertain existence profile without erasing the neighbor. Its spatial information may still be useful.

\begin{figure}[!t]
  \centering
  \includegraphics[width=\linewidth]{paper-figure2.pdf}
  \caption{Adaptive fusion-weight factorization, showing the shared quality backbone, decoupled spatial and existence-weight paths, weak structure-aware branch prior, and an FID-FIA cue applied only to the existence branch.}
  \label{fig:factorization}
\end{figure}

\begin{figure}[t]
  \centering
  \includegraphics[width=\linewidth]{paper-figure3.pdf}
  \caption{Existence-confidence weighting used in the adaptive rule: (a) the Bernoulli-level decisiveness mapping and bounded score curve, and (b) representative node-level aggregation examples showing how decisive posterior sets produce larger $\bar{c}^{(j)}$ and $q_{\mathrm{exist}}^{(j)}$.}
  \label{fig:existence-confidence}
\end{figure}

\subsection{Branch-decoupled allocation and FID-FIA refinement}

The proposed method allows spatial fusion and existence fusion to use different adaptive scores. This separation follows the Bernoulli/LMB structure because existence and spatial terms enter the fused posterior through different update blocks. Appendix~\ref{app:theory} gives the formal variational interpretation. The method therefore constructs branch-specific dedicated scores:
\begin{equation}
\label{eq:method-spatial-dedicated-score}
\tilde{\omega}_{k,s,\mathrm{sp}}^{(j)} =
m_k^{(j)}
\cdot
\left(q_{\mathrm{cov},k}^{(j)}\right)^{\alpha_x}
\cdot
\left(q_{\mathrm{link},k}^{(j)}\right)^{\beta_x},
\end{equation}
\begin{equation}
\label{eq:method-existence-dedicated-score}
\tilde{\omega}_{k,s,\mathrm{ex}}^{(j)} =
m_k^{(j)}
\cdot
\left(q_{\mathrm{link},k}^{(j)}\right)^{\beta_r}
\cdot
\left(q_{\mathrm{exist},k}^{(j)}\right)^{\alpha_r}.
\end{equation}
The powers $\alpha_x,\beta_x,\beta_r,\alpha_r \ge 0$ tune the influence of the branch-relevant cues. The final branch scores are obtained through geometric interpolation:
\begin{equation}
\label{eq:method-branch-interpolation}
\tilde{\omega}_{k,s}^{x,(j)}
=
\left(\tilde{\omega}_{k,s}^{(j)}\right)^{1-\eta_x}
\left(\tilde{\omega}_{k,s,\mathrm{sp}}^{(j)}\right)^{\eta_x},
\qquad
\tilde{\omega}_{k,s}^{r,(j)}
=
\left(\tilde{\omega}_{k,s}^{(j)}\right)^{1-\eta_r}
\left(\tilde{\omega}_{k,s,\mathrm{ex}}^{(j)}\right)^{\eta_r}.
\end{equation}
Here $\eta_x,\eta_r\in[0,1]$ are interpolation strengths: zero keeps the shared backbone, and one uses the dedicated branch score.
This separation connects one reliability backbone to two algebraically different fusion blocks. Spatial and existence consensus need not respond identically to a weight perturbation. The separation also determines how Fisher information enters the method. Cao and Zhao's construction provides a scalar indicator of target separability and posterior information \cite{CaoZhao2025InfoGeometryFusion}. In our experiments, this cue is especially effective for cardinality and existence decisions. Related multirate fusion work also supports Fisher information as a local information measure after time alignment \cite{Li2026FIMultirateAADensityFusion}. Applying one Fisher-based weight to the entire posterior can degrade spatial RMSE in the present GA-LMB setting. Appendix~\ref{app:theory} shows that this scalar baseline constrains both branches to use the same Fisher utility. Our formulation relaxes that constraint. Covariance and realized link quality continue to govern spatial fusion, whereas FID-FIA is reserved for Bernoulli existence.

\noindent\textbf{Remark 3 (why geometric interpolation).}
The shared score keeps both branches tied to the same reliability model. The dedicated spatial and existence scores then tilt the anchor toward the evidence most relevant to each branch. Geometric interpolation preserves the multiplicative reliability interpretation: a weak cue continues to constrain the final score even when another term is large.

\paragraph{Weak structure-aware refinement and stabilization.}

The evaluated Balanced and Cardinality-critical modes next apply a weak graph-aware prior after branch decoupling and before the optional FID-FIA existence modulation. This prior combines local neighborhood overlap with a soft communication-reliability modulation. For fusion node $s$ and neighbor $j\in\mathcal{N}_s$, define the neighborhood-overlap similarity
\begin{equation}
\label{eq:method-neighborhood-overlap}
\chi_s^{(j)}
=
\frac{|\mathcal{N}_s \cap \mathcal{N}_j|}
{|\mathcal{N}_s \cup \mathcal{N}_j|}.
\end{equation}
The spatial branch favors structurally redundant neighbors, while the existence branch gives only a mild novelty preference to neighbors with nonidentical neighborhoods:
\begin{equation}
\label{eq:method-structure-base-prior}
a_{x,s}^{(j)} = 1+\chi_s^{(j)},
\qquad
a_{r,s}^{(j)} = 1+\frac{1}{2}\left(1-\chi_s^{(j)}\right).
\end{equation}
Both priors are normalized to unit mean inside the local neighborhood,
\begin{equation}
\label{eq:method-structure-normalized-prior}
\xi_{x,s}^{(j)}
=
\frac{a_{x,s}^{(j)}}{|\mathcal{N}_s|^{-1}\sum_{u\in\mathcal{N}_s}a_{x,s}^{(u)}},
\qquad
\xi_{r,s}^{(j)}
=
\frac{a_{r,s}^{(j)}}{|\mathcal{N}_s|^{-1}\sum_{u\in\mathcal{N}_s}a_{r,s}^{(u)}}.
\end{equation}
When tiered packet-loss statistics are available, the implementation also forms a normalized nominal-reliability prior from the persistent drop tier of each neighbor:
\begin{equation}
\label{eq:method-structure-comm-prior}
b_s^{(j)}
=
\lambda_c+(1-\lambda_c)\left(1-p_{\mathrm{drop}}^{(j)}\right),
\qquad
\zeta_s^{(j)}
=
\frac{b_s^{(j)}}{|\mathcal{N}_s|^{-1}\sum_{u\in\mathcal{N}_s}b_s^{(u)}}.
\end{equation}
This slow prior is distinct from the realized link-quality score in the shared backbone. Specifically, $q_{\mathrm{link},k}^{(j)}$ uses delivered and dropped counts from the current step, whereas $\zeta_s^{(j)}$ summarizes the configured or long-run drop tier and does not change with the current packet outcome.
The resulting structure factors are
\begin{equation}
\label{eq:method-structure-combined-prior}
\bar{\xi}_{x,s}^{(j)}
=
\xi_{x,s}^{(j)}\left(\zeta_s^{(j)}\right)^{\delta_c},
\qquad
\bar{\xi}_{r,s}^{(j)}
=
\xi_{r,s}^{(j)}\left(\zeta_s^{(j)}\right)^{\delta_c},
\end{equation}
and they enter the two branch scores multiplicatively:
\begin{equation}
\label{eq:method-structure-modulation}
\tilde{\omega}_{k,s}^{x,(j)}
\leftarrow
\tilde{\omega}_{k,s}^{x,(j)}
\left(\bar{\xi}_{x,s}^{(j)}\right)^{\gamma_x},
\qquad
\tilde{\omega}_{k,s}^{r,(j)}
\leftarrow
\tilde{\omega}_{k,s}^{r,(j)}
\left(\bar{\xi}_{r,s}^{(j)}\right)^{\gamma_r}.
\end{equation}
Equivalently, for positive scores, the refinement adds $\gamma_x\log\bar{\xi}_{x,s}^{(j)}$ and $\gamma_r\log\bar{\xi}_{r,s}^{(j)}$ to the spatial and existence log-utilities. The main experiments use $\gamma_x=0.45$, $\gamma_r=0.08$, $\delta_c=0.30$, and $\lambda_c=0.25$, so the structure prior is deliberately weaker on the existence branch and cannot dominate the covariance, delivery, and existence-confidence backbone.

The raw adaptive weights can fluctuate sharply over time, especially when packet delivery changes abruptly. To improve stability, the normalized weights are smoothed by an exponential moving average and then passed through a minimum-weight safeguard before final renormalization. This stabilization step preserves the intended quality ordering while preventing unstable weight collapse in a time-varying communication setting.

\noindent\textbf{Remark 4 (non-Gaussian contamination and weight collapse).}
The experiments use Gaussian measurement errors, and the covariance cue comes from a moment-projected Gaussian mixture. Fallback normalization, temporal smoothing, and a final weight floor protect an active neighbor from a transient zero score. They do not make a deceptively sharp covariance trustworthy. Under impulsive or heavy-tailed contamination, a precise outlier may therefore receive excessive spatial weight. A robust implementation could instead use a shrinkage or M-estimated scatter matrix for $T_{k,i}^{(j)}$. It could also use clipped whitened residuals to form a bounded innovation factor and cap each weight before normalization. These changes preserve the separated simplex allocation, but their thresholds depend on the assumed contamination model. We do not evaluate them here. The reported results therefore provide no evidence of robustness to non-Gaussian noise.

\noindent\textbf{Remark 5 (trust and Byzantine scope).}
The present communication model treats packet loss and outages as non-adversarial and assumes that posterior parameters and delivery counters reported by active neighbors are honest. A Byzantine node could violate this assumption by advertising an artificially small covariance or a deceptive existence probability. A security-aware extension could use label-aligned spatial residuals and Bernoulli log-odds agreement to form a bounded credibility factor. Authenticated delivery metadata and per-neighbor clipping or trimming could further limit malicious influence before normalization. Resilient labeled-RFS fusion provides relevant mechanisms for this extension \cite{Gao2023ResilientLMB}. Such a layer is outside the current algorithm; the packet-loss experiments therefore establish communication robustness, not Byzantine resilience.

\paragraph{Existence-only FID-FIA refinement.}

Cardinality-critical mode adds an information-geometric refinement only to the existence branch. It accumulates pairwise Fisher information distances between locally represented targets and uses the result to modulate existence. For node $j$, the target-pair score $a_{k}^{(j)}$ is computed from measurement-updated Bernoulli components. Each component is first projected to one Gaussian. Target pairs are weighted by their existence probabilities. Their distance is then accumulated along the connecting line segment under the sensor's local Fisher metric. The metric depends on detection probability and measurement covariance. Finally, the score is normalized over the available neighborhood,
\begin{equation}
\label{eq:method-fid-normalization}
\hat{a}_{k,s}^{(j)}
=
\frac{a_{k}^{(j)}}{\epsilon + \max_{u\in\mathcal{N}_s} a_{k}^{(u)}} ,
\qquad
0 \leq \hat{a}_{k,s}^{(j)} \leq 1 .
\end{equation}
This normalized accumulation is mapped into a bounded existence-only modulation factor,
\begin{equation}
\label{eq:method-fid-cue}
q_{\mathrm{fid},k,s}^{(j)}
=
\rho_{\min}
+ (1-\rho_{\min})\left(\hat{a}_{k,s}^{(j)}\right)^{p_f}.
\end{equation}
Here $\rho_{\min}\in[0,1]$ is the lower bound of the FID-FIA existence cue and $p_f>0$ controls the nonlinearity of the normalized accumulation.
Only the existence score is modified:
\begin{equation}
\label{eq:method-fid-existence-modulation}
\tilde{\omega}_{k,s}^{r,(j)}
\leftarrow
\tilde{\omega}_{k,s}^{r,(j)}
\left(q_{\mathrm{fid},k,s}^{(j)}\right)^{\gamma_f}.
\end{equation}
Here $\gamma_f\ge0$ controls the strength of the FID-FIA modulation. The spatial score $\tilde{\omega}_{k,s}^{x,(j)}$ is left unchanged. In the main experiment, this branch uses a strong existence-only modulation and disables smoothing and flooring for the existence weights. The spatial branch retains its structure-aware stabilized weights.
Equivalently, for positive $q_{\mathrm{fid},k,s}^{(j)}$, the update adds $\gamma_f \log q_{\mathrm{fid},k,s}^{(j)}$ only to the existence-branch utility. The zero-floor setting is the boundary case that permits a zero instantaneous existence weight.

\subsection{Algorithm and operating modes}

The resulting per-node update is summarized in Algorithm~\ref{alg:adaptive-fusion}. The algorithm uses the quantities consumed by the implemented fusion layer. Spatial weights enter the Gaussian canonical-form update, whereas existence weights enter the Bernoulli existence update.

\begin{center}
\small
\fbox{%
\begin{minipage}{0.92\linewidth}
\refstepcounter{algorithm}\label{alg:adaptive-fusion}
\textbf{Algorithm~\thealgorithm. Communication-aware branch-decoupled GA-LMB fusion at node $s$ and time $k$.}
\begin{enumerate}
\setlength{\itemsep}{1pt}
\setlength{\parskip}{0pt}
\item Assemble the local neighborhood input from the delivered measurement streams and realized delivery statistics, then obtain the measurement-updated LMB posteriors $\{\pi_k^{(j)}:j\in\mathcal{N}_s\}$ for the current step.
\item Form the availability mask $m_k^{(j)}$, covariance score $q_{\mathrm{cov},k}^{(j)}$, link score $q_{\mathrm{link},k}^{(j)}$, and existence-confidence score $q_{\mathrm{exist},k}^{(j)}$.
\item Build the shared score $\tilde{\omega}_{k,s}^{(j)}$ and the spatial/existence dedicated scores $\tilde{\omega}_{k,s,\mathrm{sp}}^{(j)}$ and $\tilde{\omega}_{k,s,\mathrm{ex}}^{(j)}$.
\item Interpolate the shared and dedicated scores to obtain $\tilde{\omega}_{k,s}^{x,(j)}$ and $\tilde{\omega}_{k,s}^{r,(j)}$, then apply the weak structure-aware branch prior.
\item For Cardinality-critical mode only, compute $q_{\mathrm{fid},k,s}^{(j)}$ and multiply it into the existence score; leave the spatial score unchanged.
\item Normalize each branch over active neighbors, optionally apply temporal smoothing and a final weight floor, and write the results as $\omega_{k,s}^{x,(j)}$ and $\omega_{k,s}^{r,(j)}$.
\item Fuse each aligned Bernoulli component by using $\omega_{k,s}^{x,(j)}$ in the Gaussian canonical-form update and $\omega_{k,s}^{r,(j)}$ in the Bernoulli existence update.
\end{enumerate}
\end{minipage}}
\end{center}

\paragraph{Computational cost.}
Let $A_{k,s}=|\mathcal{A}_{k,s}|$ be the number of active neighbors and $M_k$ the number of aligned Bernoulli labels. Let $\bar G_k$ be the mean mixture size per local Bernoulli component and $d_x$ the state dimension. Forming the covariance backbone requires one moment projection per active neighbor and label, for $O(A_{k,s}M_k\bar G_k d_x^2)$ time. The remaining weight operations add $O(A_{k,s}M_k+A_{k,s})$ time. Neighborhood-overlap priors are computed once with the local graph, and their per-update modulation costs $O(A_{k,s})$. Thus Balanced-mode weight formation is linear in both active neighbors and aligned labels for fixed state dimension and mixture size.

For each aligned Bernoulli component, the projected-Gaussian GA merger costs $O\!\left(A_{k,s}(\bar G_k d_x^2+d_x^3)\right)$. The first term is moment projection; the second covers covariance solves, inverses, and determinants. The Bernoulli existence product adds only $O(A_{k,s})$. Cardinality-critical mode additionally evaluates all retained target pairs. With $Q$ line-segment quadrature points and $C_F$ the cost of one fixed-dimension Fisher-metric evaluation, its FID-FIA term costs $O(A_{k,s}M_k^2 Q C_F)$ in the worst case. The implementation uses $Q=3$. This quadratic pair term, rather than branch normalization, explains the larger measured runtime of FID-FIA-based modes in Section~\ref{sec:results}.

\paragraph{Operating modes.}

The method family supports two operating modes. \emph{Balanced mode} uses the three-factor backbone, separate spatial and existence scores, a weak graph prior, and temporal stabilization. It leaves the final FID-FIA modulation inactive. This mode is intended for deployments constrained by localization accuracy, RMSE stability, runtime, or energy.

\emph{Cardinality-critical mode} adds FID-FIA only to the existence branch. It is intended for settings where target-number errors, missed tracks, or false support outweigh moderate costs in RMSE and runtime. Both modes share the same backbone and branch algebra. They differ only in the final existence cue. The experiments show the resulting tradeoff. Balanced mode favors localization at low cost, whereas Cardinality-critical mode favors cardinality agreement with substantially more computation.

\paragraph{Optional extension modules.}

The evaluated software also supports several exploratory extension modules. Appendix~\ref{app:additional-modules} summarizes them because their empirical effects are smaller, more coupled, or less stable than the evidence for the main three-factor backbone plus the retained branch-specific refinements.
